% ============================================================
%  Part I — Pronouns
% ============================================================
\gpart{II}{Pronouns}

% ------------------------------------------------------------
\gsec{Personal pronouns}

\gintro{A personal pronoun stands in for a noun, and like the noun it replaces
it changes shape according to its job in the sentence. Formal \textit{Sie}
appears twice below, because one set of forms covers both a single person and a
group.}

\begin{casetbl}{}
  Number & Person & English & Nominativ & Akkusativ & Dativ & Genitiv \\
  \SetCell[r=6]{c,m} Singular & 1st & I              & ich & mich & mir   & meiner \\
                              & 2nd & you \eng{inf.}   & du  & dich & dir   & deiner \\
                              & 2nd & you \eng{formal} & Sie & Sie  & Ihnen & Ihrer  \\
                              & 3rd & he             & er  & ihn  & ihm   & seiner \\
                              & 3rd & she            & sie & sie  & ihr   & ihrer  \\
                              & 3rd & it             & es  & es   & ihm   & seiner \\
  \SetCell[r=4]{c,m} Plural   & 1st & we             & wir & uns  & uns   & unser  \\
                              & 2nd & you \eng{inf.}   & ihr & euch & euch  & euer   \\
                              & 2nd & you \eng{formal} & Sie & Sie  & Ihnen & Ihrer  \\
                              & 3rd & they           & sie & sie  & ihnen & ihrer  \\
\end{casetbl}

\begin{gnote}
Formal \textit{Sie} has one set of forms for both numbers, so the two rows above
are identical. It always takes plural verb endings — \textit{Sie sind} — whether
you are addressing one person or twenty.
\end{gnote}

\begin{gnote}
The genitive column is given for completeness. Genitive personal pronouns are
archaic and survive mainly in fixed phrases such as \textit{Wir gedenken ihrer};
ordinary speech uses \textit{von + Dativ} instead. \textit{unserer} and
\textit{eurer} are also found.
\end{gnote}

\tcap{Four different words spelled \textit{ihr}}

\begin{flattbl}{colspec={X[0.7,l] X[1,l] X[2.3,l]}}
  ihr & nominative & \emph{you} (plural, informal) — \textit{ihr seid} \\
  ihr & dative     & \emph{to her} — \textit{ich gebe ihr das Buch} \\
  ihr & possessive & \emph{her} or \emph{their} — \textit{ihr Hund} \\
  Ihr & possessive & \emph{your} (formal) — \textit{Ihr Hund} \\
\end{flattbl}

% ------------------------------------------------------------
\gsec{Reflexive pronouns}

\gintro{A reflexive pronoun is used when the subject and the object are the same
person, as in \textit{ich wasche \textbf{mich}} — \emph{I wash myself}. Many
German verbs are reflexive where their English equivalents are not.}

\begin{gtbl}{colspec={X[0.75,c,m] X[0.95,l] X[1,l] X[0.95,c] X[0.85,c]}, hline{3-Y}={0.25pt,solid,Line}}
  Number & Person & Subject & Akkusativ & Dativ \\
  \SetCell[r=4]{c,m} Singular & 1st              & ich & mich & mir  \\
                              & 2nd \eng{inf.}   & du  & dich & dir  \\
                              & 2nd \eng{formal} & Sie & sich & sich \\
                              & 3rd              & er / sie / es & sich & sich \\
  \SetCell[r=4]{c,m} Plural   & 1st              & wir & uns  & uns  \\
                              & 2nd \eng{inf.}   & ihr & euch & euch \\
                              & 2nd \eng{formal} & Sie & sich & sich \\
                              & 3rd              & sie & sich & sich \\
\end{gtbl}

\begin{gnote}
Only the first and second person singular distinguish accusative from dative
(\textit{mich}/\textit{mir}, \textit{dich}/\textit{dir}). Every other row is the
same in both cases, so in practice there are two pairs to keep apart and
\textit{sich} covers the rest.
\end{gnote}

\tcap{Which case does the reflexive take?}

\begin{flattbl}{colspec={X[1,l] X[2.3,l]}}
  Akkusativ & when there is no other object: \textit{Ich wasche \textbf{mich}.} \\
  Dativ     & when a direct object is already present: \textit{Ich wasche \textbf{mir} die Hände.} \\
\end{flattbl}

\tcap{Common reflexive verbs}

\begin{flattbl}{colspec={X[1.15,l] X[1,l] X[1.15,l] X[1,l]}}
  sich freuen        & to be glad       & sich setzen      & to sit down \\
  sich erinnern      & to remember      & sich beeilen     & to hurry \\
  sich interessieren & to be interested & sich fühlen      & to feel \\
  sich vorstellen    & to imagine \eng{dat.} & sich anziehen & to get dressed \\
  sich befinden      & to be located    & sich entscheiden & to decide \\
\end{flattbl}

% ------------------------------------------------------------
\gsec{Relative pronouns}

\gintro{A relative pronoun introduces a clause describing a noun already
mentioned: \emph{the man \textbf{who} lives here}. In German it looks almost
exactly like the definite article, and it sends the verb of its clause to the
end.}

\begin{gendertbl}{}
  Case & maskulin & feminin & neutrum & Plural \\
  Nominativ & der & die & das & die \\
  Akkusativ & den & die & das & die \\
  Dativ     & dem & der & dem & \en{denen} \\
  Genitiv   & \en{dessen} & \en{deren} & \en{dessen} & \en{deren} \\
\end{gendertbl}

\begin{gnote}
Five cells differ from the definite article: \textit{denen} in the dative
plural, and all four genitive forms. The remaining eleven are identical.
\end{gnote}

\begin{gnote}
Two independent decisions produce the right form. The \textbf{gender} comes from
the noun being described; the \textbf{case} comes from the job the pronoun does
inside its own clause. Neither one constrains the other.
\end{gnote}

\tcap{Worked examples}

\begin{extbl}{colspec={X[1.15,l] X[1.85,l]}}
  Der Mann, \textbf{der} dort steht      & subject $\rightarrow$ nominative, masc. \\
  Der Mann, \textbf{den} ich kenne       & direct object $\rightarrow$ accusative, masc. \\
  Der Mann, \textbf{dem} ich helfe       & \textit{helfen} governs the dative \\
  Die Frau, \textbf{deren} Auto neu ist  & possession $\rightarrow$ genitive, fem. \\
  Die Leute, mit \textbf{denen} ich rede & \textit{mit} governs the dative plural \\
\end{extbl}

\begin{gnote}
The clause is fenced off with commas and its verb sits at the end. After
\textit{alles, nichts, etwas, vieles}, and after a whole preceding clause, the
relative pronoun is \textit{was} rather than a der-form.
\end{gnote}

% ------------------------------------------------------------
\gsec{Indefinite pronouns}

\gintro{Indefinite pronouns refer to people or things without identifying them:
\emph{someone}, \emph{nothing}, \emph{all}. Some decline for case, others never
change at all.}

\tcap{Referring to people}

\begin{gtbl}{colspec={X[0.8,l] X[0.8,c] X[0.9,c] X[0.9,c] X[1.1,l]}}
  Word & Nom. & Akk. & Dat. & Meaning \\
  man     & man     & einen     & einem     & one, you, people \\
  jemand  & jemand  & jemanden  & jemandem  & someone \\
  niemand & niemand & niemanden & niemandem & no one \\
  jeder   & jeder   & jeden     & jedem     & everyone \\
\end{gtbl}

\begin{gnote}
\textit{man} exists only in the nominative and borrows from \textit{einer} in
the other cases. It is the usual German alternative to the passive:
\textit{\textbf{Man} spricht hier Deutsch} where English says \emph{German is
spoken here}. It is unrelated to \textit{der Mann} and is never capitalised.
\end{gnote}

\tcap{Referring to things}

\begin{flattbl}{colspec={X[0.8,l] X[1,l] X[0.8,l] X[1,l]}}
  etwas   & something     & nichts  & nothing \\
  alles   & everything    & vieles  & much, many things \\
  einiges & a fair amount & manches & some things \\
\end{flattbl}

\tcap{Quantifiers used on their own}

\begin{flattbl}{colspec={X[0.8,l] X[1,l] X[0.8,l] X[1,l]}}
  alle    & all     & einige & some, a few \\
  viele   & many    & wenige & few \\
  mehrere & several & beide  & both \\
  keiner  & none    & andere & others \\
\end{flattbl}

\begin{gnote}
An adjective after \textit{etwas} or \textit{nichts} is capitalised, treated as
neuter and takes the strong ending: \textit{etwas Gut\textbf{es}},
\textit{nichts Neu\textbf{es}}. The prefix \textit{irgend-} adds vagueness:
\textit{irgendjemand} \emph{somebody or other}, \textit{irgendwo}
\emph{somewhere or other}, \textit{irgendwann} \emph{at some point}.
\end{gnote}
